% from ega4-2-fr.tex
\oldpage[IV-2]{172}
L'assertion (i) a déjà été vue au cours de la démonstration de II) et est
indépendante de l'hypothèse dimensionnelle sur $A$. L'hypothèse que le
morphisme $\operatorname{Spec}(A')\to\operatorname{Spec}(A)$ est normal
entraîne que $A'\otimes_A k=A'/\mathfrak mA'$ est un anneau normal local et
noethérien, donc intègre, ce qui montre déjà que
$\mathfrak p'=\mathfrak mA'$ est premier. On a, d'après (6.1.2),
$\dim(A'_{\mathfrak p'})=\dim(A)+
\dim(A'_{\mathfrak p'}/\mathfrak mA'_{\mathfrak p'})$. Mais
$A'_{\mathfrak p'}/\mathfrak mA'_{\mathfrak p'}=
A'_{\mathfrak p'}/\mathfrak p'A'_{\mathfrak p'}$ est le corps résiduel de
$A'_{\mathfrak p'}$, donc $\dim(A'_{\mathfrak p'})=1$. Le fait que
$A'_{\mathfrak p'}$ est contenu dans le composé direct des $L'_j$ tels que
$\mathfrak q'_j\subset\mathfrak p'$ résulte de ce que les idéaux
$\mathfrak q'_jA'_{\mathfrak p'}$ pour ces indices $j$ sont les idéaux
premiers minimaux de $A'_{\mathfrak p'}$.

Reste à prouver (6.14.1.3). On a évidemment
$A'\subset K'\cap\overline{A'_{\mathfrak p'}}$. Inversement, soit $y'$ un
élément de cette intersection; soit $a$ un «paramètre» pour $A$, donc tel
que $Aa$ soit $\mathfrak m$-primaire, et soit $a'$ l'image de $a$ dans
$A'$; tout élément de $K$ peut s'écrire $x/a^n$ pour $x\in A$ et un $n>0$
puisque $Aa$ contient une puissance de $\mathfrak m$; donc on peut écrire
$y'=x'/a'^n$ avec $x'\in A'$. Notons maintenant que $\mathfrak p'$ est le
seul idéal premier associé à
$A'/a'^nA'=(A/a^nA)\otimes_A A'$: comme $A'$ est un $A$-module plat, cela
résulte de (3.3.1), $\mathfrak m$ étant le seul idéal premier associé à
$A/a^nA$ et $k\otimes_A A'$ étant intègre. Par suite, $a'^nA'$ est l'image
réciproque dans $A'$ de $a'^nA'_{\mathfrak p'}$; comme par hypothèse
$y'\in\overline{A'_{\mathfrak p'}}$, l'image de $x'$ dans
$A'_{\mathfrak p'}$ appartient à $a'^nA'_{\mathfrak p'}$; d'où
$x'\in a'^nA'$ et $y'\in A'$, ce qui achève la démonstration.

Ce lemme étant établi, dans le cas où nous sommes ramenés à la fin de II),
$B'$ est radiciel sur $A'$ (I, 3.5.7) et par suite est aussi un anneau
local; en outre $B'$ est entier sur $A'$, donc, si $\mathfrak q'$ est
l'unique idéal premier de $B'$ au-dessus de $\mathfrak p'$,
$\mathfrak q'B'_{\mathfrak p'}$ est le seul idéal maximal de
$B'_{\mathfrak p'}$ (Bourbaki, Alg. comm., chap. V, § 2, n$^o$ 1, prop.
1), donc $B'_{\mathfrak q'}=B'_{\mathfrak p'}=B\otimes_A
A'_{\mathfrak p'}$. Nous allons d'abord démontrer que
$B'_{\mathfrak p'}$ est un anneau noethérien et normal. Or, comme $B$
contient $A$ et est contenu dans $K$ et que $A'_{\mathfrak p'}$ est un
$A$-module plat, $B'_{\mathfrak p'}$ contient $A'_{\mathfrak p'}$ et est
contenu dans $K'_{\mathfrak p'}$, donc dans le produit $L_1$ des $L'_j$
tels que $\mathfrak q'_j\subset\mathfrak p'$. Pour tout indice $j$ tel que
$\mathfrak q'_j\subset\mathfrak p'$, soit $\mathfrak a'_j$ le produit des
$L'_h$ tels que $h\neq j$, de sorte que
$\mathfrak a'_j\cap A'=\mathfrak q'_j$; comme tout élément de
$A'-\mathfrak p'$ est régulier dans $L_1$, on a aussi
\[
\mathfrak a'_j\cap A'_{\mathfrak p'}=
\mathfrak q'_jA'_{\mathfrak p'}.
\]
Soit $\mathfrak r'_j=B'_{\mathfrak p'}\cap\mathfrak a'_j$, de sorte que
$B'_{\mathfrak p'}/\mathfrak r'_j$ est isomorphe à la projection de
$B'_{\mathfrak p'}$ dans $L'_j$; donc $B'_{\mathfrak p'}/\mathfrak r'_j$
contient l'anneau local intègre de dimension 1,
$A'_{\mathfrak p'}/\mathfrak q'_jA'_{\mathfrak p'}$, et est contenu dans
son corps des fractions $L'_j$; il est par suite noethérien en vertu du
théorème de Krull-Akizuki (Bourbaki, Alg. comm., chap. VII, § 2, n$^o$ 5,
prop. 5).

